
\begin{center}
	\Large \textbf{Preface}
\end{center}
\vspace{0.5cm}
The \gls*{DVIC} is a community of makers that develops technologies
within philosophical and critical frameworks to shape our societies’
futures. The objective is to implement real-world solutions as well as
design projects to enhance public engagement, improve education, and
overall provide scientific knowledge. Our researchers contribute actively
to top-level international research in multiple fields, including artificial
intelligence, human-computer interactions, education, and ecology. We
believe that these objectives require a transdisciplinary approach, that
bridges the gap between sciences, techniques, sociology, and philosophy.
This is performed by collaborating with other scientists and industrial and
startup sharing our values, to form strong research partnerships...\\
\\
The Artificial Lives group, led by Dr. Clement Duhart, aims to develop
the next generation of machines and Human-Machine Interfaces.
The group members strongly believe that through the combination of
Design and Engineering, human-centered technologies can blend into our
environments to become invisible, vastly improving daily lives. To achieve
this vision, the members contribute to human-computer interactions,
cognitive enhancement through new forms of extended intelligence,
learning platforms, and cobotic. Our bio-inspired, multidisciplinary
approach couples AI and virtual reality with intelligent materials, robotics
and the Internet of Things.\\
\\
For the past two years, De Vinci Innovation Center (DVIC) students
following the Creative Technologies curriculum had the opportunity to
develop their vision on technology, innovation, and society.
This proceeding is a composition of six master’s theses, ranging
from Machine Learning, Human-Computer-Interaction to Robotics. The
authors strongly believe that developing alternative futures requires
new types of engineering that take into consideration both the people’s
needs and the environment. These documents have been written to reflect
this vision and refined over several months with an iterative reviewing
supervised by the Principal Investigators.\\
\\
The Authors, the Principal Investigators and the whole DVIC community
is proud of releasing this first proceeding. We dedicate this first edition to Pascal Brouaye and Nelly Rouyres, without whom nothing would have been possible.
\begin{figure}[ht!]
	\centering
	\includegraphics[height=1.5cm]{./images/dvic.png}
	\hspace{1cm}
	\includegraphics[height=1.5cm]{./images/AL.png}
	\hspace{1cm}
	\includegraphics[height=1.5cm]{./images/pulv.png}
\end{figure}

\newpage
\paragraph{ }
\newpage
